\section{RL Through a Bounded, Evicting Cache}
\label{sec:method}

\draft{Method section. Pull details from \texttt{keys\_values/rl/grpo/} and
\texttt{docs/grpo\_integration.md}. Aim for one figure (the step diagram) and
one algorithm block.}

\subsection{Overview}
Each RL step has three phases, all executed through the same bounded cache of
$K$ slots:
\begin{enumerate}
  \item \textbf{Chunked prefill.} The prompt is fed in chunks of $C$ tokens
    (we use $C=1024$). After each chunk the cache scores entries and evicts
    down to $K$.
  \item \textbf{Generation.} $G$ completions are sampled from the evicting
    cache; per-token log-probabilities are recorded and reused as
    $\log\pi_{\text{old}}$.
  \item \textbf{Chunked backward.} The loss is back-propagated chunk by chunk
    with bounded temporary memory, re-materialising the cache state that was
    live for each chunk rather than holding the full sequence of activations.
\end{enumerate}

\subsection{Shared-prompt prefill}
With $G$ completions per prompt the prefill is shared: the prompt is
processed once and the cache state is branched per completion
(response-only schedule, cf.\ \draft{LongStraw cite}). This is token-exact
relative to the na\"ive schedule and gives a $1.5\times$ generation speedup at
7k-token prompts.

\subsection{Cache configuration}
\draft{Describe H2O scoring, 8-bit KV quantisation (q8), the \texttt{cl}
budget, and why full-coverage caches OOM in the gradient pass (memory is
paid in backward too).} The configurations used in the paper are listed in
\cref{tab:cache_configs}.

\begin{table}[t]
  \caption{Cache configurations. \draft{Fill from experiment configs.}}
  \label{tab:cache_configs}
  \centering\small
  \begin{tabular}{lccc}
    \toprule
    Name & Policy & KV dtype & Budget $K$ \\
    \midrule
    dense-default   & none & bf16 & $=$ seq.\ len.\\
    h2o-q8 (8k campaign) & H2O & int8 & \draft{TODO}\\
    h2o-q8 (32k flagship) & H2O & int8 & 8192\\
    h2o-q8 (wide)   & H2O & int8 & 12288\\
    \bottomrule
  \end{tabular}
\end{table}

\subsection{Memory accounting}
\label{sec:memory_accounting}
\draft{Short derivation: dense peak grows as $O(L)$ in KV plus $O(L)$
activations per chunk; bounded cache peak is $O(K)$ plus per-chunk terms, so
flat in $L$. Tie to \cref{tab:memory_frontier}.}

\subsection{Optimiser and scale}
Full fine-tuning of 7B models on 48\,GB cards uses an 8-bit paged optimiser
\citep{dettmers2022optimizers}. \draft{Mention LoRA r=16 was tried and did
not learn (memory note) only if we want to pre-empt the reviewer question.}
